\chapter{Proof of the main theorem}\label{sec:b7}

In this section, we prove \cref{thm:main}. By \eqref{eq:in-envelope},
four times the fractional-moment error $e_{N,\beta}(\lambda)$ is the
integral of the defect $\calD_N(\rho)$ over $\rho\in(0,1]$. We split
this integral at a cutoff $\rho_0$. Above the cutoff, we verify the
stability assumption of the cavity argument with
\cref{prop:lz-stability} and bound the defect with
\cref{cor:c7-defect}. Below the cutoff, we bound the defect in \cref{prop:sm-small}, using the
quantile-shell estimate of \cref{prop:wi-vfs} and the root localization
of \cref{cor:lz-omega}. The cutoff is $\rho_0=\min\{1,K\eta^{1/7}\}$,
where $K$ depends only on $\beta$ and
$\eta=e_{N,\beta}(\lambda)+L_s\lambda^5$ is the trial budget. This budget
contains the unknown error, and the two estimates combine
into an implicit inequality for $e_{N,\beta}(\lambda)$. In
\cref{subsec:b7-general}, we choose the tilt $\lambda=N^{-7/50}$,
normalize this inequality by $N^{-7/10}$, and close it. Then
\cref{prop:fnd-mean-bias} converts the bound on
$e_{N,\beta}(N^{-7/50})$ into the bound on the mean bias. In
\cref{rem:b7-exponents}, we explain how the exponent $7/10$ arises.

We use results from \cref{sec:fnd} (the defect identity
\eqref{eq:in-envelope}, the bound \eqref{eq:fnd-moving}, the
trial budget of \cref{lem:fnd-trial}, the trace bound \eqref{eq:in-Lvar},
the decomposition of \cref{lem:fnd-decomp}, and the reduction of the
mean bias in \cref{prop:fnd-mean-bias}), from \cref{sec:wi,sec:th}
(\cref{prop:wi-vfs,prop:th-j}), from \cref{sec:lz-root}
(\cref{cor:lz-omega}), from \cref{sec:sc} (\cref{prop:lz-stability}), and
from \cref{sec:c7-closure} (\cref{cor:c7-defect}). Before the proof of
\cref{thm:main}, we list these inputs with the parameters at which they
are used.

\paragraph{Standing assumptions and notation.} Throughout this section,
$\beta>1$ is fixed, and $C_\beta$ denotes a constant depending only on
$\beta$ whose value may change from line to line. The constants
$\lambda_0\le\min\{1,\beta/4\}$ and $L_s\ge1$ are those of
\cref{sec:fnd}; in particular, \eqref{eq:fnd-scalar-trial-cost} holds
for $0<\lambda\le2\lambda_0$. Recall that $c=\beta^2\rho$, $M=Nc$,
$s=\lambda/\beta$, and $\mu=s\delta_0+\nu$, and that
$\log_+x=\log(\max\{1,x\})$ for $x\ge0$. As in \cref{sec:fnd}, a
\emph{trial budget at $\rho$} is a number $\eta\ge0$ that bounds the left
side of \eqref{eq:fnd-trial}. Since every minimizer $\gamma_*$ of
$F_\rho$ attains $\min_\gamma F_\rho(\gamma)$, such an $\eta$ is a trial
budget for every minimizer at $\rho$ in the sense of
\eqref{eq:loc-trialbudget}.

\section{The interval near \texorpdfstring{$\rho=0$}{rho = 0}}
\label{subsec:b7-small}

We estimate the three terms of $\calD_N(\rho)=\beta(A+2J+V)$ separately.
By the trace bound \eqref{eq:in-Lvar} and
\eqref{eq:th-J}, we have at every minimizer
\begin{equation}\label{eq:sm-AJ}
 A\le\frac1{Ms},\qquad
 J\le\frac{341(1+\log_+M)^2}{Ms}.
\end{equation}
To estimate the self-overlap variance $V$, we combine
\cref{prop:wi-vfs} on dyadic quantile shells with the root localization
of \cref{cor:lz-omega}. This follows Steps~2--5 of the proof of
\cite[Proposition~7.7]{AL}, with \cref{prop:wi-vfs} in place of
\cite[Lemma~3.3]{AL} and \cref{cor:lz-omega} in place of
\cite[Lemma~7.6]{AL}.

\begin{proposition}[The interval near $\rho=0$]\label{prop:sm-small}
Let $N\ge2$, $0<\lambda\le\lambda_0$, and $0<\rho\le1$, and let $\eta>0$
be a trial budget at $\rho$, that is, a bound for the left side of
\eqref{eq:fnd-trial}. Let
$\omega(\rho,\eta)=\min\{1,(\eta/\rho)^{1/6}\}$ be as in
\eqref{eq:lz-omega}. Then every minimizing profile $\mu$ at these
parameters satisfies
\begin{align}
 \beta(A+2J+V)
 &\le C_\beta\left[
 \frac{\log^2(2+N)}{N\rho\lambda}
 +(N\rho)^{-2/3}
 +(N\rho)^{-2/3}\lambda^{-1/3}\omega(\rho,\eta)^{2/3}
 \right],\label{eq:sm-pointwise}\\
 \beta(A+2J+V)&\le C_\beta(N\rho)^{-1/2},
 \label{eq:sm-moving}
\end{align}
where $A$, $J$, and $V$ are computed at $\mu$. In particular, the same
bounds hold for $\calD_N(\rho)$ whenever $\calD_N(\rho)$ is defined. If
$0<\rho_0\le1$ and $\eta$ is a trial budget at every $\rho\in(0,\rho_0]$,
then
\begin{equation}\label{eq:sm-integrated}
 \int_0^{\rho_0}\calD_N(\rho)\dd\rho
 \le C_\beta\left[
 N^{-1}+\frac{\log^3(2+N)}{N\lambda}
 +N^{-2/3}\rho_0^{1/3}
 +N^{-2/3}\lambda^{-1/3}\eta^{1/9}\rho_0^{2/9}
 \right].
\end{equation}
The constants depend only on $\beta$; in particular, they do not depend
on $N$, $\rho$, $\lambda$, $\eta$, the minimizer, or the positions and
masses of the atoms of $\nu$.
\end{proposition}

\begin{proof}
The proof has five steps. We partition the quantile variable into dyadic
shells and bound the diameter of the image of each shell by the root
localization. We then apply \cref{prop:wi-vfs} on each shell and sum
over the shells to bound $V$. Finally, we deduce the pointwise bounds
and integrate them in $\rho$.

\paragraph{Step 1: Shells and localization.} Fix a minimizing profile
$\mu=s\delta_0+\nu$, and let $Q$ be the quantile function of $\nu$ from
\eqref{eq:fnd-quantile}. Recall that $Q$ pushes Lebesgue measure on
$(0,1-s)$ forward to $\nu$ and that $a(Q(v))\ge s+v$ for
$0<v<1-s$. In particular,
\begin{equation}\label{eq:sm-V-quantile}
 V=\int\Var(S_q)\,\nu(\dd q)=\int_0^{1-s}\Var(S_{Q(v)})\dd v .
\end{equation}
For the dyadic levels $u=2^js<1$, $j\ge0$, of which there are finitely
many, we consider the shells
\[
 I_u=\{v\in(0,1-s):u\le s+v<2u\},\qquad h_u=|I_u|.
\]
(Here $u$ denotes a level; it is unrelated to the parameter of the
cavity interpolation in \cref{sec:c7}.) For each $v\in(0,1-s)$, we have
$s<s+v<1$, and there exists exactly one $j\ge0$ such that
$2^js\le s+v<2^{j+1}s$; for this $j$, we have $u=2^js<1$. Hence the
shells partition $(0,1-s)$. Each $I_u$ is an interval of positive length
$h_u\le u$, and $a(Q(v))\ge s+v\ge u$ for $v\in I_u$. In other words,
$I_u$ has threshold $u$ in the sense of \cref{sec:wi-setup}. A shell may contain only
part of the quantile interval of an atom of $\nu$, which
\cref{prop:wi-vfs} allows.

We write $\omega=\omega(\rho,\eta)$, and we let
$\ell(I_u)=\operatorname*{ess\,sup}_{v,w\in I_u}|Q(v)-Q(w)|$ be the
diameter of the image of $I_u$, as in \cref{sec:wi-interval}. Since
$\mu$ is a minimizer, $0<\lambda\le\lambda_0$, and $\eta$ is a trial
budget for $\gamma_*=\beta\mu([0,\cdot\,])$, \cref{cor:lz-omega} implies
that $0\le Q(v)\le C_Q(s+v+\omega)$ for $0<v<1-s$, with $C_Q$ depending
only on $\beta$. For $v\in I_u$, we have $s+v<2u$. As $Q\ge0$, this
gives
\[
 \ell(I_u)\le C_Q(2u+\omega)\le2C_Q(u+\omega).
\]

\paragraph{Step 2: The variance on one shell.} We use \cref{prop:wi-vfs}
on $I_u$, with threshold $u$, length $h_u$, diameter $\ell(I_u)$, and
$\epsilon_u=(Mu)^{-1}$. The first term in \eqref{eq:wi-vfs} is at most
$\epsilon_u(1+\log_+(h_u/\epsilon_u))\le(Mu)^{-1}(1+\log_+(Mu^2))$,
because $h_u/\epsilon_u=Muh_u\le Mu^2$. For the second term, we use
$h_u\le u$ and $(u+\omega)^{2/3}\le u^{2/3}+\omega^{2/3}$ and find
\[
\begin{aligned}
 (\epsilon_u\ell(I_u))^{2/3}h_u^{1/3}
 &\le(2C_Q)^{2/3}(Mu)^{-2/3}(u+\omega)^{2/3}u^{1/3}\\
 &\le(2C_Q)^{2/3}M^{-2/3}
       \bigl(u^{1/3}+u^{-1/3}\omega^{2/3}\bigr).
\end{aligned}
\]
We conclude that
\begin{equation}\label{eq:sm-shell}
 \int_{I_u}\Var(S_{Q(v)})\dd v
 \le C_\beta\left[
 \frac{1+\log_+(Mu^2)}{Mu}
 +M^{-2/3}\bigl(u^{1/3}+u^{-1/3}\omega^{2/3}\bigr)
 \right].
\end{equation}
The term $u^{-1/3}\omega^{2/3}$ accounts for the range of levels that
the localization permits beyond the shell's own mass.

\paragraph{Step 3: Summation over the shells.} For $u=2^js$, we have
$\log_+(Mu^2)\le\log_+(Ms^2)+2j\log2$. Using
$\sum_{j\ge0}2^{-j}=\sum_{j\ge0}j2^{-j}=2$, we see that the first terms
of \eqref{eq:sm-shell} sum to at most
\[
 \frac{2(1+\log_+(Ms^2))+4\log2}{Ms}
 \le\frac{5(1+\log_+(Ms^2))}{Ms},
\]
where we used $4\log2<3$. Let $c_3=(1-2^{-1/3})^{-1}$. Summing downward
from the largest level, which is less than one, we have
$\sum_{u<1}u^{1/3}\le c_3$, and summing upward from $s$, we have
$\sum_{u<1}u^{-1/3}\le c_3s^{-1/3}$. As the shells partition
$(0,1-s)$, summing \eqref{eq:sm-shell} and using
\eqref{eq:sm-V-quantile} yields
\begin{equation}\label{eq:sm-V}
 V\le C_\beta\left[
 \frac{1+\log_+(Ms^2)}{Ms}
 +M^{-2/3}\bigl(1+s^{-1/3}\omega^{2/3}\bigr)
 \right].
\end{equation}

\paragraph{Step 4: The pointwise bounds.} We insert \eqref{eq:sm-AJ} and
\eqref{eq:sm-V} into $\beta(A+2J+V)$ and use
\[
 Ms=\beta N\rho\lambda,\qquad
 \beta M^{-2/3}=\beta^{-1/3}(N\rho)^{-2/3},\qquad
 \beta M^{-2/3}s^{-1/3}=(N\rho)^{-2/3}\lambda^{-1/3}.
\]
As $M\le\beta^2N$ and $Ms^2=N\rho\lambda^2\le N$, both $1+\log_+M$ and
$1+\log_+(Ms^2)$ are at most $C_\beta\log(2+N)$. Using also
$\log(2+N)\ge1$, we find that $\beta A$, $2\beta J$, and $\beta$ times the
first term of \eqref{eq:sm-V} are at most
$C_\beta\log^2(2+N)/(N\rho\lambda)$, while the other two terms of
\eqref{eq:sm-V} give the last two terms of \eqref{eq:sm-pointwise}. This
proves \eqref{eq:sm-pointwise}. No lower bound on $M$ is needed here.

For \eqref{eq:sm-moving}, recall the first bound in
\eqref{eq:fnd-moving}: at every minimizer,
$\int\E(R_q-q)^2\,\nu(\dd q)\le C\min\{1,M^{-1/2}\}$ with a numerical
constant $C$. The measure $\nu$ is carried by the contact levels, where
\eqref{eq:in-decomp} gives $\E(R_q-q)^2=a^{\rm th}_q+2j_q+\Var(S_q)$.
Integrating against $\nu$, we find
$\beta(A+2J+V)=\beta\int\E(R_q-q)^2\,\nu(\dd q)\le C\beta M^{-1/2}
=C(N\rho)^{-1/2}$, which is \eqref{eq:sm-moving}.

At every $\rho$ at which $\calD_N(\rho)$ is defined, a minimizer
exists, and $\calD_N(\rho)=\beta(A+2J+V)$ for every minimizer at $\rho$
by \eqref{eq:in-envelope}. This implies that the two bounds also hold
for $\calD_N(\rho)$.

\paragraph{Step 5: Integration.} Suppose that $\eta$ is a trial budget
at every $\rho\in(0,\rho_0]$. By Step~4, the bounds \eqref{eq:sm-pointwise}
and \eqref{eq:sm-moving} hold for $\calD_N(\rho)$ at almost every
$\rho\in(0,\rho_0]$. As they hold for every minimizer, no measurable
choice of minimizers is needed to integrate them. We first use
\eqref{eq:sm-moving}, whose singularity at $\rho=0$ is integrable:
\[
 \int_0^{\min\{\rho_0,1/N\}}\calD_N(\rho)\dd\rho
 \le C_\beta N^{-1/2}\int_0^{1/N}\rho^{-1/2}\dd\rho
 =2C_\beta N^{-1}.
\]
If $\rho_0\le1/N$, this proves \eqref{eq:sm-integrated}. If
$\rho_0>1/N$, we integrate \eqref{eq:sm-pointwise} over $[1/N,\rho_0]$.
The first term contributes at most
\[
 \frac{\log^2(2+N)}{N\lambda}\int_{1/N}^{\rho_0}\frac{\dd\rho}\rho
 =\frac{\log^2(2+N)\log(N\rho_0)}{N\lambda}
 \le\frac{\log^3(2+N)}{N\lambda},
\]
and the second at most
$N^{-2/3}\int_0^{\rho_0}\rho^{-2/3}\dd\rho=3N^{-2/3}\rho_0^{1/3}$. For the
third term, we use
$\omega(\rho,\eta)^{2/3}=\min\{1,(\eta/\rho)^{1/9}\}\le(\eta/\rho)^{1/9}$,
which holds for all $\rho>0$. With this bound, the third term
contributes at most
\[
 N^{-2/3}\lambda^{-1/3}\eta^{1/9}
       \int_0^{\rho_0}\rho^{-2/3-1/9}\dd\rho
 =\frac92N^{-2/3}\lambda^{-1/3}\eta^{1/9}\rho_0^{2/9}.
\]
Adding the contributions, we obtain \eqref{eq:sm-integrated}. The
constants in Steps~1--5 depend only on $\beta$, through $C_Q$, through
the absolute constants of \cref{prop:wi-vfs}, \cref{prop:th-j}, and
\eqref{eq:fnd-moving}, and through powers of $\beta$.
\end{proof}

\section{Closure at the critical tilt}
\label{subsec:b7-general}\label{subsec:b7-main}

We now fix the tilt. We take $\lambda=N^{-7/50}$, so that the scalar
cost $\lambda^5$ in the trial budget equals $N^{-7/10}$. The cutoff
$\rho_0$ is defined through the trial budget, which contains the
unknown error $e_{N,\beta}(\lambda)$, and we do not assume any
preliminary bound on this error. After division by $N^{-7/10}$, the
terms that depend on the budget become sublinear powers of the single
variable $x=\eta/N^{-7/10}$, and the resulting inequality bounds $x$.
We follow the closure in the proof of \cite[Proposition~7.7]{AL}:
there, the trial budget \cite[(7.44)]{AL} also contains the unknown
error, and the resulting inequality \cite[(7.52)]{AL} is closed in
Step~6 because the error enters only through sublinear powers. Here the
cutoff produces the leading term $N^{-2/3}\eta^{1/21}$.

\paragraph{The inputs.} The proof of \cref{thm:main} uses the following
results, at the parameters listed. In each case, the hypotheses are
verified in the step indicated.

\begin{center}
\small
\begin{tabular}{@{}>{\raggedright\arraybackslash}p{0.2\textwidth}
 >{\raggedright\arraybackslash}p{0.4\textwidth}
 >{\raggedright\arraybackslash}p{0.32\textwidth}@{}}
\toprule
Result & Conclusion used & Parameters in the proof\\
\midrule
\Cref{lem:fnd-trial} &
$\eta=e+L_s\lambda^5$ is a trial budget at every $\rho\in(0,1]$ &
$N\ge N_1$, $\lambda=N^{-7/50}\le\lambda_0$ (Step~1)\\
\eqref{eq:in-envelope} &
$4e=\int_0^1\calD_N(\rho)\dd\rho$, and $\calD_N(\rho)=\beta(A+2J+V)$ at
every minimizer, for almost every $\rho$ &
$s=\lambda/\beta\le1/4$ (Steps~1 and~4)\\
\Cref{prop:lz-stability} &
$J^{\scal}_{a,d}\preceq(1+\eps_N)I$ with
$\eps_N\le C_{\rm st}[\rho^{-3/4}\eta^{1/4}+\lambda^2+N^{-1}]$, at every
minimizer &
$N\ge2$, $\rho\in[\rho_0,1]$, trial budget $\eta$ (Step~3)\\
\Cref{cor:c7-defect}(c) &
$\int_{\rho_0}^1\calD_N\le C_\beta\log^2(2+N)/(N\lambda\rho_0)$ &
$1/N\le\rho_0\le1$, $s\le1/2$, stability on $[\rho_0,1]$ (Step~3)\\
\Cref{prop:sm-small} &
the integrated bound \eqref{eq:sm-integrated} &
$0<\rho_0\le1$, trial budget $\eta>0$ on $(0,\rho_0]$ (Step~4)\\
\Cref{prop:fnd-mean-bias} &
$(P_\beta-p_{N,\beta})/\beta\le32L_s\lambda^5+2e$ &
$N\ge N_1$, $0<\lambda\le\lambda_0$ (Step~6)\\
\bottomrule
\end{tabular}
\end{center}

\begin{proof}[Proof of \cref{thm:main}]
The proof has six steps. We fix the parameters and the common trial
budget, define the cutoff, bound the integral of the defect above and
below the cutoff, close the resulting inequality after normalization,
and convert the bound on the fractional-moment error into the bound on
the mean bias.

\paragraph{Step 1: Parameters and the trial budget.} Let
$N_1=\max\{2,\lceil\lambda_0^{-50/7}\rceil\}$, which depends only on
$\beta$. For every integer $N\ge N_1$, we have $N^{-7/50}\le\lambda_0$,
and we set
\begin{equation}\label{eq:b7-parameters}
 \lambda=N^{-7/50},\qquad r=N^{-7/10}=\lambda^5,\qquad
 e=e_{N,\beta}(\lambda),\qquad \eta=e+L_sr.
\end{equation}
Then $0<\lambda\le\lambda_0\le\min\{1,\beta/4\}$ and
$s=\lambda/\beta\le1/4$. By \cref{lem:fnd-trial}, $\eta$ is a trial
budget at every $\rho\in(0,1]$. The number $e$ is finite. Indeed, the
constrained value $E_\beta(\lambda)$ is finite, and
$F^c_{N,\beta}/\beta$ is a Lipschitz function of finitely many Gaussian
variables (\cref{sec:fnd}), which has finite exponential moments of
every order; these moments enter \eqref{eq:fnd-error-definition}. As
$e\ge0$ by \eqref{eq:in-envelope} and $L_s\ge1$, we have
\begin{equation}\label{eq:b7-eta-lower}
 \infty>\eta\ge L_s\lambda^5\ge\lambda^5>0.
\end{equation}

\paragraph{Step 2: The cutoff.} Let $C_{\rm st}\ge1$ be the constant of
\cref{prop:lz-stability}. We define
\begin{equation}\label{eq:b7-cutoff}
 K=12C_{\rm st},\qquad \rho_0=\min\{1,K\eta^{1/7}\}.
\end{equation}
In both cases $\rho_0<1$ and $\rho_0=1$, we have
\[
 \rho_0\le K\eta^{1/7},\qquad
 \rho_0\ge\min\{1,K\lambda^{5/7}\}\ge\lambda^{5/7}=N^{-1/10}\ge N^{-1},
\]
where we used \eqref{eq:b7-eta-lower}, $K\ge1$, and $\lambda\le1$.

\paragraph{Step 3: The stable interval.} Suppose first that $\rho_0<1$,
and let
$\rho\in[\rho_0,1]$. Then $\rho\ge K\eta^{1/7}\ge K\lambda^{5/7}$. Since
$N^{-1}\le\lambda\le1$, this implies
\[
 \lambda^2\le\lambda^{5/7}\le\frac\rho K,\qquad
 N^{-1}\le\lambda^{5/7}\le\frac\rho K,\qquad
 \rho^{-3/4}\eta^{1/4}\le\frac\rho{K^{7/4}}\le\frac\rho K,
\]
where the last bound follows from $\rho^{7/4}\ge K^{7/4}\eta^{1/4}$ and
$K\ge1$. Let $\mu$ be any minimizer at $\rho$, with cumulative function
$a$. The hypotheses of \cref{prop:lz-stability} hold: $N\ge2$,
$0<\rho\le1$, $0<\lambda\le\lambda_0$, and $\eta$ is a trial budget for
$\mu$. Applying the proposition with $d=\beta^2(1-t/N)$, the cavity
coefficient of \eqref{eq:cavityparameters}, we have
\[
 J^{\scal}_{a,d}\preceq(1+\eps_N)I\ \text{ on }\mathcal H_\nu,\qquad
 0\le\eps_N\le C_{\rm st}\bigl[\rho^{-3/4}\eta^{1/4}+\lambda^2+N^{-1}\bigr]
 \le\frac{3C_{\rm st}\rho}K=\frac\rho4,
\]
where $\eps_N\ge0$ is the number defined in \cref{prop:lz-stability}.
This is the stability condition \eqref{eq:c7-stab}, at every minimizer
and every $\rho\in[\rho_0,1]$. As $N\ge2$, $1/N\le\rho_0\le1$, and
$s\le1/2$, \cref{cor:c7-defect}(c) applies. Using
$\rho_0=K\eta^{1/7}\ge\lambda^{5/7}$ and
$\lambda^{12/7}=N^{-6/25}$, we obtain
\begin{equation}\label{eq:b7-closure}
 \int_{\rho_0}^1\calD_N(\rho)\dd\rho
 \le C_\beta\frac{\log^2(2+N)}{N\lambda\rho_0}
 \le C_\beta\frac{\log^2(2+N)}{N\lambda^{12/7}}
 =C_\beta N^{-19/25}\log^2(2+N).
\end{equation}
When $\rho_0=1$, the integral is zero, so \eqref{eq:b7-closure} holds in
both cases.

\paragraph{Step 4: The interval near zero.} The hypotheses of
\cref{prop:sm-small} hold on $(0,\rho_0]$: we have $N\ge2$,
$0<\lambda\le\lambda_0$, and $0<\rho_0\le1$. By Step~1 and
\eqref{eq:b7-eta-lower}, $\eta>0$ is a trial budget at every
$\rho\in(0,\rho_0]$. By Step~2, we have $\rho_0\le K\eta^{1/7}$ in both
cases $\rho_0<1$ and $\rho_0=1$, which implies that
\[
 \rho_0^{1/3}\le K^{1/3}\eta^{1/21},\qquad
 \eta^{1/9}\rho_0^{2/9}\le K^{2/9}\eta^{1/9+2/63}=K^{2/9}\eta^{1/7}.
\]
We insert these bounds into \eqref{eq:sm-integrated}, add
\eqref{eq:b7-closure}, and use $4e=\int_0^1\calD_N(\rho)\dd\rho$ from
\eqref{eq:in-envelope}. This gives
\[
\begin{aligned}
 e\le C_\beta\Bigl[
 &N^{-1}+N^{-43/50}\log^3(2+N)+N^{-19/25}\log^2(2+N)\\
 &+N^{-2/3}\eta^{1/21}+N^{-2/3}\lambda^{-1/3}\eta^{1/7}\Bigr],
\end{aligned}
\]
where we used $N^{-1}\lambda^{-1}=N^{-43/50}$. As $1>43/50>19/25$ and
$\log(2+N)\ge1$, the first three terms are at most
$3N^{-19/25}\log^3(2+N)$. We conclude that
\begin{equation}\label{eq:b7-normalized-input}
 e\le C_\beta\left[
 N^{-19/25}\log^3(2+N)+N^{-2/3}\eta^{1/21}
                     +N^{-2/3}\lambda^{-1/3}\eta^{1/7}\right].
\end{equation}

\paragraph{Step 5: Normalization and closure.} Let
\[
 x=\frac\eta r=\frac er+L_s ,
\]
which is finite by \eqref{eq:b7-eta-lower} and satisfies $x\ge L_s\ge1$.
Then $\eta=rx$, and the terms of the bound for $e$ in Step~4 have the
following sizes after division by $r=N^{-7/10}$.
\begin{center}
\small
\begin{tabular}{@{}lll@{}}
\toprule
Term & Origin & Size after division by $r$\\
\midrule
$N^{-1}$ & $\rho\le1/N$, by \eqref{eq:sm-moving} & $N^{-3/10}$\\
$N^{-43/50}\log^3(2+N)$ & first term of \eqref{eq:sm-pointwise} &
$N^{-4/25}\log^3(2+N)$\\
$N^{-19/25}\log^2(2+N)$ & stable interval, \eqref{eq:b7-closure} &
$N^{-3/50}\log^2(2+N)$\\
$N^{-2/3}\eta^{1/21}$ & second term of \eqref{eq:sm-pointwise} & $x^{1/21}$\\
$N^{-2/3}\lambda^{-1/3}\eta^{1/7}$ & third term of \eqref{eq:sm-pointwise} & $N^{-1/50}x^{1/7}$\\
\bottomrule
\end{tabular}
\end{center}
The entries in the last column are exact. For the two terms that depend
on the budget, we have
\[
 N^{-2/3}\eta^{1/21}=N^{-2/3}r^{1/21}x^{1/21}
 =N^{-2/3-1/30}x^{1/21}=rx^{1/21}
\]
and
\[
 N^{-2/3}\lambda^{-1/3}\eta^{1/7}
 =N^{-2/3+7/150-1/10}x^{1/7}=N^{-18/25}x^{1/7}=rN^{-1/50}x^{1/7},
\]
since $2/3+1/30=7/10$ and $18/25=7/10+1/50$. For the first term of
\eqref{eq:b7-normalized-input}, we have $N^{-19/25}/r=N^{-3/50}$.
Dividing \eqref{eq:b7-normalized-input} by $r$ and adding $L_s$ to both
sides, we find
\begin{equation}\label{eq:b7-normalized}
 x\le L_s+C_\beta\left[
 N^{-3/50}\log^3(2+N)+x^{1/21}+N^{-1/50}x^{1/7}\right]
 \le C_\beta x^{1/7}.
\end{equation}
In the last inequality, we used that
$\sup_{N\ge2}N^{-3/50}\log^3(2+N)$ is finite, because $\log(2+N)$ grows
more slowly than any power of $N$. We also used $N^{-1/50}\le1$ and
the bounds $1\le x^{1/7}$ and $x^{1/21}\le x^{1/7}$, which hold because
$x\ge1$.
Dividing \eqref{eq:b7-normalized} by $x^{1/7}>0$, we find
$x^{6/7}\le C_\beta$, that is, $x\le C_\beta^{7/6}$. Then
\[
 e\le\eta=rx\le C_\beta N^{-7/10}.
\]
This argument uses no preliminary bound on $e$ and no smallness of
$\eta$; it uses only that $x$ is finite.

\paragraph{Step 6: The mean bias.} By \cref{prop:fnd-mean-bias}, which
applies since $0<\lambda\le\lambda_0$, and by Step~5, we have
\[
 0\le\frac{P_\beta-p_{N,\beta}}\beta
 \le32L_s\lambda^5+2e
 \le(32L_s+2C_\beta)N^{-7/10}.
\]
Multiplying by $\beta$, we obtain
$P_\beta-p_{N,\beta}\le C_\beta N^{-7/10}$ for every $N\ge N_1$. For the
finitely many integers $2\le N<N_1$, the mean bias $P_\beta-p_{N,\beta}$
is finite. The number
\[
 C^{\rm small}_\beta=\max_{2\le N<N_1}N^{7/10}(P_\beta-p_{N,\beta}),
\]
with the maximum interpreted as zero when $N_1=2$, is therefore finite
and depends only on $\beta$. Replacing $C_\beta$ by $\max\{C_\beta,C^{\rm small}_\beta\}$, we
obtain the bound of \cref{thm:main} for every $N\ge2$. This completes the
proof.
\end{proof}

\begin{remark}\label{rem:b7-exponents}
The leading term in \eqref{eq:b7-normalized-input} is
$N^{-2/3}\eta^{1/21}$. The exponent $1/21$ is the product of the exponent
$1/3$ of $\rho_0$ in the integrated shell estimate
\eqref{eq:sm-integrated} and the exponent $1/7$ of $\eta$ in the cutoff
\eqref{eq:b7-cutoff}, which comes from the stability threshold
$\rho^{-3/4}\eta^{1/4}\lesssim\rho$. If $\eta$ is of order $N^{-a}$,
the leading term is of the same order when $a=2/3+a/21$, that is,
$20a/21=2/3$ and $a=7/10$. The scalar cost $\lambda^5$ in the trial
budget is then of the same order when $\lambda=N^{-7/50}$. For this
choice, the localization correction has the extra factor $N^{-1/50}$
after normalization, and the stable interval and the two elementary
terms in \eqref{eq:sm-integrated} contribute at most
$N^{-3/50}\log^3(2+N)$, as the table in Step~5 shows. The bound
\eqref{eq:b7-closure} on the stable interval has only a squared
logarithm; the third logarithm comes from integrating the first term of
\eqref{eq:sm-pointwise} over $[1/N,\rho_0]$. These contributions have
strict polynomial slack, and the leading term carries no logarithm.
After the closure, $L_sN^{-7/10}\le\eta\le C_\beta N^{-7/10}$, so that
for all sufficiently large $N$ the cutoff satisfies $\rho_0<1$ and is
comparable to $N^{-1/10}$.
\end{remark}
